\chapter{Imagens, tabelas e matemática}
\label{chap:elementos}

\section{Figuras}

O pacote \arquivo{graphicx} fornece \comando{includegraphics}. Coloque imagens
em uma pasta própria e use o ambiente \arquivo{figure} para legenda, numeração e
referência:

\begin{lstlisting}[language=LaTeXBook]
\begin{figure}[htbp]
  \centering
  \includegraphics[width=0.75\textwidth]{images/diagrama.pdf}
  \caption{Fluxo de producao do livro.}
  \label{fig:fluxo}
\end{figure}
\end{lstlisting}

As letras \arquivo{htbp} indicam posições preferidas: aqui, topo, base ou página
de figuras. São sugestões, pois LaTeX escolhe uma posição que preserve a boa
composição da página. Use PDF para desenhos vetoriais e PNG ou JPEG para imagens
fotográficas. Evite capturas de tela de texto quando o conteúdo puder ser
escrito diretamente.

\section{Tabelas}

O ambiente \arquivo{table} fornece legenda; \arquivo{tabular} organiza células.
O pacote \arquivo{booktabs} cria linhas com proporções profissionais:

\begin{lstlisting}[language=LaTeXBook]
\begin{table}[htbp]
  \centering
  \caption{Classes comuns de documento.}
  \label{tab:classes}
  \begin{tabular}{ll}
    \toprule
    Classe & Uso principal \\
    \midrule
    article & textos curtos \\
    report  & relatorios longos \\
    book    & livros com capitulos \\
    \bottomrule
  \end{tabular}
\end{table}
\end{lstlisting}

Em \verb|{ll}|, cada letra representa uma coluna alinhada à esquerda. Use
\arquivo{c} para centralizar e \arquivo{r} para alinhar à direita. Não use linhas
verticais sem necessidade: espaço e alinhamento geralmente separam melhor os
dados.

\section{Expressões matemáticas}

Matemática curta pode aparecer no parágrafo, como $E=mc^2$. Expressões que
merecem destaque usam um ambiente próprio:

\begin{lstlisting}[language=LaTeXBook]
\begin{equation}
  a^2 + b^2 = c^2
  \label{eq:pitagoras}
\end{equation}
\end{lstlisting}

O pacote \arquivo{amsmath} acrescenta ambientes para alinhar várias expressões:

\begin{lstlisting}[language=LaTeXBook]
\begin{align}
  x + y &= 10 \\
  x - y &= 2
\end{align}
\end{lstlisting}

O símbolo \verb|&| marca o ponto de alinhamento. Não use modo matemático apenas
para obter itálico; isso altera espaçamento e significado dos caracteres.

\begin{pratica}
  Adicione ao seu documento uma imagem com legenda, uma tabela de três linhas e
  uma equação numerada. Crie rótulos e cite os três elementos no texto com
  \comando{ref}.
\end{pratica}

